\documentclass[UTF8]{ctexbook}

\title{中学英语语法手册}

\author{王乐知\and徐晨曦}

\begin{document}

\maketitle

\begin{center}
{\heiti{\LARGE 前言}}
\end{center}

本书使用\LaTeX{}技术排版，系作者临时起意，本书的诞生纯属偶然，内容错误请向作者指出。

本书是英语教材的总结和补充，因此很多地方较为简略（也有作者的技术原因）。

本书不是教材，请不要把本书当成教材来学习，至多作为教材的梳理和补充。

在接受了这些前提后，请继续浏览。

\tableofcontents

\part{简单句和基本语法}

\chapter{单词的分类、变化和搭配}

{\heiti{单词（word）是构成短语（phrase），句子（sentence）和文章（article）的基础，因此在学习语法之前，应当先掌握单词的相关知识。}}

\section{单词的分类}

根据词性，单词可以分为如下几类：

名词（Noun n.）指代人、事物、地点或抽象概念。

动词（Verb v.）指代动作或状态。

形容词（Adjective adj.）指代用来修饰名词的特征。

副词（Adverb adv.）指代修饰动词/形容词/其他副词，表程度、方式等。

代词（Pronoun pron.）指代替代名词的词。

介词（Preposition prep.）指代表示方位、时间、方式等关系的词。

连词（Conjunction conj.）指代连接单词、短语或句子的词。

感叹词（Interjection int.）指代表达情感或惊呼的词。

冠词（Article art. ）指代限定名词的特指或泛指。

数词（Numeral num.）指代表示数量或顺序的词。

下面让我们分别了解一下它们。

\subsection{名词}

在中学，名词（n.）可以被简略地分为可数名词（Countable Noun cn.）和不可数名词（Uncountable Noun un.），其中可数名词指代可以用数目数清的名词，不可数名词指代无法用数目数清的名词。

\subsubsection{可数名词}

可数名词有单复数的变化。单数指一个该物体，复数指多个该物体。因此名词常与数词，冠词等限制词（如a/an/the/my等）连用。在复数时无需限制词。

e.g. There is an apple. 有一个苹果。/There are two apples. 有两个苹果。

可数名词的变化规则见后。

\subsubsection{不可数名词}

不可数名词没有单复数的变形，则无需限定词的修饰。

e.g. Water is polluted. 水被污染。

被动语态见后。

\subsection{动词}

动词可以被分成如下的几类：

及物动词（Transitive Verb vt.）指代后面必须接宾语的动词。

不及物动词（Intransitive Verb vi.）指代后面不需要接宾语的动词。

系动词（Linking Verb link-v.）指代连接主语和表语，表状态或特征的动词。

助动词（Auxiliary Verb aux.）指代帮助构成时态、语态或疑问的动词。

情态动词（Modal Verb modal.）指代表能力、可能、必须等语气的动词。

\subsubsection{及物动词}

及物动词本身不可以表达完整的含义，需要借助后边的宾语才能表达完整含义。

e.g.I ate an apple. 我吃了一个苹果。

动词的变化规则见后。

\subsubsection{不及物动词}

不及物动词本身可以表达完整的含义，故而不需要宾语，如果要接宾语，则可以通过加上介词后接宾语，其与介词和及物动词的关系可以写成如下的公式：

\[
vi.+prep.=vt.
\]

e.g. I am eating.我正在吃。/ I ate an apple.我吃了一个苹果。 

\subsubsection{系动词}

系动词（Linking Verbs）是英语动词中最特殊的一类：它不表示动作，而是说明主语的属性。在中学中可以被分成这几类：

1. 状态系动词，表示“存在”或“持续的状态”：be

2. 感官系动词，表示“感官感知”：look, smell, taste, sound, feel	

3. 变化系动词，表示“从一种状态变为另一种”：become, get, grow, turn, go	

4. 保持系动词，表示“维持某种状态不变”：keep, remain, stay	

5. 表象系动词，表示“看起来像”，但未必是事实：seem, appear

以上可以简记为：一be,两似乎，三保持，四变，五感官，再加上一个go。

\subsubsection{助动词}

助动词没有实际含义，帮助实义动词构成时态、语态、疑问句、否定句或强调语气。常见的助动词有do,have,be等。

\subsection{形容词}

形容词没有详细的分类，但是从含义角度来说，可以分为品质、大小、颜色、国籍、材料等。

\subsection{副词}

副词可以分为方式、程度、时间、频率、地点、原因副词等。

\subsection{代词}

代词可以分为人称、物主、反身、指示、不定、疑问、关系代词等。

\subsection{介词}

介词可以分为时间、地点、方式、原因/目的介词等。

\subsection{连词}

连词可以分为并列连词和从属连词。

\subsection{感叹词}

感叹词从感情上可以被分为惊喜/赞叹、痛苦/烦恼、犹豫/停顿、呼唤/应答等。

\subsection{冠词}

冠词分为不定冠词、定冠词、零冠词。

\subsection{数词}

数词可以分为基数词和序数词。

\section{单词的变化}

在上一节中，我们简要介绍了单词的分类，这一节我们来学习单词的变化。

注：介绍次序按照1.1的单词词性顺序。

\subsection{名词}



\end{document}